% ============================================================================
% Paper 34 Sprint 2 Track 8 draft: Wick rotation / signature change
%
% Target slot: §III.27 (between gauge-choice §III.26 and apparatus identity
%               §III.28 in Sprint 2's planned numbering)
% Label: sec:proj_wick_rotation
%
% §IV three-axis dictionary row (for PM integration into the §IV table):
%   Projection   : Wick rotation / signature change
%   Variables    : signature (binary: Riemannian Euclidean / Lorentzian
%                  Minkowskian); rapidity parameter not introduced here
%                  (full Lorentz-boost projection at the operator-system
%                  level is the open extension)
%   Dimension    : preserves (analytic continuation; the physical observable
%                  carries the same dimension on both sides, only the
%                  interpretation changes)
%   Transcendental: $2\pi \cdot \mathbb{Q}$ via $\mathrm{Vol}(S^1)$ on the
%                  Euclidean time circle (M1 Hopf-base measure mechanism;
%                  same generator as §III.15 observation/temporal-window,
%                  structurally distinct projection: analytic continuation
%                  vs Euclidean compactification)
%   Layer        : Layer 2 → Layer 2 (operates on the already-projected
%                  Euclidean spectral triple, not on the bare graph)
%
% Cross-references used:
%   - sec:proj_observation (sibling: same M1 generator, different mechanism)
%   - sec:proj_restmass (Dirac single-particle γ factor distinct from
%                        signature-change content)
%   - sec:proj_spectral_action (M1 mechanism on heat-kernel trace)
%   - sec:layer2_boundary (Layer-2 projection class)
%   - Paper 32 §VIII Remark rem:bisognano_wichmann_reading (text reference)
%   - Paper 35 §VIII subsec:bisognano_wichmann_reading (text reference)
%   - Paper 35 §VIII subsec:unruh_pendant (text reference)
%
% §VIII open-question REWORDING RECOMMENDED for PM integration:
% The current §VIII open-question entry "Lorentz boost as a candidate
% twenty-sixth projection" (~lines 5335-5510 of paper_34_projection_taxonomy.tex)
% should be reworded to reflect that the metric-functional-level Wick rotation
% has been promoted to §III.27, and that the open question is now the full
% operator-system-level Lorentzian extension (Lorentzian propinquity,
% modular Hamiltonian computation at finite n_max, full Lorentz-boost
% projection between rapidities). The numbering "twenty-sixth" becomes
% structurally awkward post-promotion; the entry should be renamed something
% like "Operator-system-level Lorentzian extension" and recast as the open
% deepening of §III.27 rather than as an entirely separate candidate. The
% Bisognano-Wichmann reading content from that §VIII entry is now subsumed
% by §III.27's "Empirical anchor" and "Honest scope" paragraphs.
% ============================================================================

\subsection{Wick rotation / signature change}
\label{sec:proj_wick_rotation}
\textbf{Source:} the framework's bare spectral-triple machinery on
the Fock-projected $S^3$ graph (or any natural-geometry Euclidean
spectral triple, including the $T_{S^3} \otimes T_{S^1_\beta}$
tensor-product construction of Sprint TD Track~1).  By construction
every operator in the GeoVac inventory ---
Camporesi--Higuchi Dirac on round $S^3$, Bargmann--Segal heat
kernel on $S^5$, Wilson plaquette functionals on Hopf graphs, all
$\zeta_R(2k)$ values entering through $\zeta_{D^2}$ on closed
Riemannian manifolds --- is computed in Euclidean (Riemannian)
signature.\\
\textbf{Target:} the Lorentzian (Minkowski-signature) physical
observable obtained by analytic continuation
$\tau \to i t$, $\beta \to i \beta$ on the imaginary-time circle
(or proper-time circle in the Rindler case).  Concrete instances
under the Wick-rotation chain
[Hartle--Hawking 1976]~$\to$~[Sewell 1982]~$\to$~[Bisognano--Wichmann
1976]~$\to$~[Unruh 1976]:
the Hartle--Hawking KMS state at
$\beta_{\mathrm{cigar}} = 8\pi M$ on the Euclidean Schwarzschild
cigar continues to the Hawking thermal state on the Lorentzian
Schwarzschild exterior at $T_H = 1/(8\pi M)$; the periodic
Euclidean Rindler coordinate $\eta_E$ with period $2\pi$ continues
to the Lorentzian Rindler boost at acceleration $a$, with the
accelerated-observer Unruh temperature $T_U = a/(2\pi)$; the modular
automorphism of the algebra of observables localized in a Lorentzian
Rindler wedge, with respect to the Minkowski vacuum, is identified
by Bisognano--Wichmann with the unitary representation of the Lorentz
boost subgroup preserving the wedge, and the analytic period of the
modular flow is the same $2\pi = \mathrm{Vol}(S^1)$ that appears on
the Euclidean side.\\
\textbf{Variables introduced:} the signature index --- a discrete
binary choice between Riemannian Euclidean and Lorentzian
Minkowskian.  Structurally analogous to the gauge-choice projection
(Sprint~2 companion track): a discrete bookkeeping selection that
fixes how the framework's spectral content is read as a physical
observable.  Continuous rapidity parameters relating two Lorentzian
frames at relative velocity $v/c$ are \emph{not} introduced by this
projection; a full Lorentz-boost projection at the operator-system
level (acting on all particles simultaneously, relating two truncated
spectral triples at different observer rapidities) is a separate
structural extension named in \S~\ref{sec:open}.\\
\textbf{Dimension introduced:} dimension-preserving.  Analytic
continuation does not introduce a new $[L]$, $[E]$, $[M]$, or $[T]$
scale; it changes the interpretation of an already-dimensioned
quantity from a Euclidean (imaginary-time-periodic) to a Lorentzian
(real-time-thermal or wedge-modular) reading.  In the V/D/P tagging
of \S~\ref{sec:dictionary}, the projection sits in the same
dimension-preserving class as Sturmian reparameterization
(\S~\ref{sec:proj_sturmian}) and the gauge-choice projection
(Sprint~2 companion track).\\
\textbf{Transcendental signature:} $2\pi \cdot \mathbb{Q}$ via
$\mathrm{Vol}(S^1)$ on the Euclidean time circle (Hopf-base measure /
M1 sub-mechanism of the master Mellin engine; see Paper~32~§VIII
Remark on the Bisognano--Wichmann reading and Paper~35~§VIII for the
framework-level structural correspondence).  Critically:\ this is
the \emph{same generator} as the observation / temporal-window
projection (\S~\ref{sec:proj_observation}), which injects
$2\pi \cdot \mathbb{Q}$ per Matsubara mode via compactification of
Euclidean time at finite $\beta = 1/T$.  The two projections share
the M1 transcendental signature but are \emph{structurally distinct}
as mechanisms: observation / temporal-window compactifies a single
Euclidean time direction at finite period $\beta$ and reads out
finite-temperature observables (Casimir at $T \neq 0$,
Stefan--Boltzmann free energy, Matsubara loops), whereas Wick
rotation analytically continues between two signatures and reads out
Lorentzian observables (Hawking radiation, Unruh thermal bath,
Bisognano--Wichmann modular flow) from the same Euclidean data.
Both surface the M1 $2\pi$ on the same $S^1$ factor;
the structural reading of what that $2\pi$ \emph{means} differs:\
a Hopf-bundle measure factor on the Riemannian side; a
modular-flow / Lorentz-boost orbit period on the Lorentzian side
via the published Wick-rotation prescription.\\
\textbf{Used in:} the Hawking-temperature off-precision catalogue
row in \S~\ref{sec:offprecision} (error class C, $M$ external
Einstein-equations input, machinery-witness only); Sprint TD Track~4's
reproduction of $T_H = 1/(8\pi M)$ on the Euclidean Schwarzschild
cigar from M1 alone; Sprint Unruh-pendant's reproduction of
$T_U = a/(2\pi)$ on the Euclidean Rindler wedge from the same M1
mechanism; the structural-correspondence reading of
$\beta = 2\pi / (\text{surface gravity})$ that unifies Hawking,
Sewell, Bisognano--Wichmann, and Unruh as four faces of a single
Wick-rotation theorem.  Any future framework-side observable whose
physical interpretation requires Lorentzian signature (Wightman-class
expectation value of a real-time-ordered operator product, KMS state
on a Lorentzian wedge algebra, real-time Schwinger--Keldysh closed
time contour) consumes this projection.\\
\textbf{Empirical anchor:} the four-witness Wick-rotation theorem,
each face reproducing a shared $2\pi = \mathrm{Vol}(S^1)$ from the
M1 mechanism:\
(i)~Hawking temperature $T_H = 1/(8\pi M) = \kappa_g / (2\pi)$ on the
Euclidean Schwarzschild cigar, with surface gravity
$\kappa_g = 1/(4M)$ (Hartle--Hawking~1976~\cite{hartle_hawking1976};
Sprint~TD Track~4 reproduces this from M1 with sympy residual zero
on all three identities $\beta = 8\pi M$, $T_H = \kappa_g/(2\pi)$,
$\beta = 4M \cdot 2\pi$);
(ii)~Sewell's identification of the Hartle--Hawking thermal state
restricted to the Schwarzschild exterior as a KMS state at
$\beta = 8\pi M$ on a globally hyperbolic Lorentzian manifold with
a bifurcate Killing horizon~\cite{sewell1982};
(iii)~Bisognano--Wichmann identification of the modular automorphism
of the algebra of observables localized in a Rindler wedge, with
respect to the vacuum state, with the unitary representation of the
Lorentz boost subgroup preserving the
wedge~\cite{bisognano_wichmann1976}, with analytic period
$\sigma_{i \cdot 2\pi} = $ identity forced by Tomita--Takesaki KMS;
(iv)~Unruh temperature $T_U = a / (2\pi)$ for a uniformly accelerated
observer at proper acceleration $a$, derived in the same Wick-rotation
construction with $\beta_{\mathrm{Rindler}} = 2\pi / a$ on the
Euclidean Rindler $\eta_E$-circle~\cite{unruh1976} (Sprint
Unruh-pendant reproduces this verbatim with the same M1 mechanism,
sympy residual zero, lowest bosonic Matsubara mode $= a$, lowest
fermionic mode $= a/2 = \pi T_U$).
All four faces reproduce the shared $2\pi$ from the framework's
M1 Hopf-base measure on the same $S^1$ factor; the published
Wick-rotation prescription supplies the bridge to the Lorentzian-side
modular-flow / boost-orbit vocabulary in each case.\\
\textbf{Honest scope.}  This projection commits to a
\emph{structural correspondence}, not a literal identification.
The framework demonstrably reproduces $\beta = 2\pi / \kappa_g$ on
all four Wick-rotation faces from M1 alone, and the published
Hartle--Hawking / Sewell / Bisognano--Wichmann / Unruh chain
identifies that $2\pi$ with the Lorentzian-side modular-flow /
Lorentz-boost orbit period.  The identification stands on the
published QFT machinery, which is mature and independently verified.
The framework does \emph{not} provide an operator-level proof
internal to the truncated spectral triple
$\mathcal{T}_{n_{\max}}$:\ the corresponding statement
(compute the modular Hamiltonian $K$ on $\mathcal{T}_{n_{\max}}$
for a Rindler-wedge-restricted Camporesi--Higuchi state, verify
$\sigma_{i\cdot 2\pi} = $ identity at finite $n_{\max}$, and take
the Gromov--Hausdorff limit using the lemmas of Paper~38) is the
principal falsifier of the structural correspondence, named with
estimated cost 4--8 weeks and MEDIUM-HIGH priority in Paper~32~§VIII
Remark on the Bisognano--Wichmann reading.

Two further scope distinctions are required.
First, \emph{Wick rotation is not the full Lorentzian extension.}
This projection lifts the framework's metric-functional Riemannian
output (a Mellin trace, a heat-kernel coefficient, a Wilson
expectation value) to its Lorentzian analytic continuation in
cases where the bridge is a single Wick rotation of a temporal
circle.  It does \emph{not} construct a Lorentzian (Minkowski-signature)
spectral triple with Krein-space Hilbert space and fundamental
symmetry $\mathcal{J}^2 = I$; that extension would carry the
framework's signature index from $(m,n) = (3, 0)$ in the
Bizi--Brouder--Besnard~2018~\cite{bizi_brouder_besnard2018}
classification to $(3, 1)$.  A clean Lorentz-boost projection at
the metric / curvature level, acting on \emph{all} particles in the
system simultaneously and relating two truncated metric spectral
triples at different observer rapidities, requires that full
Lorentzian-signature spectral triple plus a Lorentzian analog of
the Latr\'emoli\`ere propinquity used in Paper~38~\cite{paper38}
for the WH1 PROVEN result.  The propinquity literature is purely
Riemannian as of May~2026 (Latr\'emoli\`ere, Toyota,
Hekkelman--McDonald, Leimbach--van~Suijlekom); constructing a
Lorentzian propinquity is original NCG-mathematics work at the
6--12 month scale.  That open extension is flagged in
\S~\ref{sec:open} and is the structural deepening of this
\S~\ref{sec:proj_wick_rotation} entry, not a separate candidate.

Second, \emph{distinct from existing $\gamma$-extensions in the
catalogue.}  The Dirac-sector $\gamma = \sqrt{1 - (Z\alpha)^2}$ in
the Camporesi--Higuchi spinor lift (\S~\ref{sec:proj_spinor}) is a
bound-state radial Lorentz factor for a single-particle hydrogenic
wavefunction in a central potential; it is not a frame
transformation between two observers, and the Wick rotation
discussed here does not introduce or modify it.  The rest-mass
projection (\S~\ref{sec:proj_restmass}) rescales a single particle's
mass under $m_e \to m_\text{red}$ within Euclidean signature; the
Breit retardation projection
(\S~\ref{sec:proj_breit_retardation}) introduces the rest-mass
ratio $m_l/m_n$ as a two-body kinematic control parameter.  None of
these is a signature change.  Wick rotation sits structurally above
all three:\ those projections operate within a fixed Euclidean
spectral triple; this projection toggles between the Euclidean and
Lorentzian readings of the \emph{same} computed quantity, with the
toggle bridged by the published Wick-rotation chain.\\
\textbf{Structural reading.}  Wick rotation sits in the
Layer-2~$\to$~Layer-2 class (cf.\
\S~\ref{sec:layer2_boundary}):\ it operates on already-projected
Euclidean spectral content, not on the bare graph or its initial
projection chain.  Structurally it is a
\emph{metric-functional-level} projection:\ the framework computes a
Mellin trace, heat-kernel coefficient, or other Euclidean
spectral-action functional (Paper~32~§VIII case-exhaustion theorem
ensures the transcendental content is M1, M2, or M3); Wick rotation
then reads that functional as a Lorentzian physical observable via
the published Hartle--Hawking / Sewell / Bisognano--Wichmann / Unruh
chain.

The structural-correspondence-not-literal-identification verdict
crystallizes the framework's position on Lorentzian physics in
general:\ at the metric functional level (specifically the M1
$2\pi = \mathrm{Vol}(S^1)$ signature on a Wick-rotated time circle),
the framework's Euclidean output IS the Lorentzian modular-flow /
boost-orbit period, with no extension required.  At the
operator-system level (Tomita--Takesaki modular automorphism acting
on the truncated algebra at finite $n_{\max}$, Krein-space
representation, frame-level Lorentz-boost between truncated metric
spectral triples at different rapidities), the framework currently
sees the same $2\pi$ but cannot internally close the period
identity; that closure is the named open extension.

Three observations distinguish this projection from
\S~\ref{sec:proj_observation} despite shared M1 generator and
shared $2\pi$ transcendental content.
First, mechanism:\ observation / temporal-window compactifies a
single Euclidean time direction at finite period $\beta = 1/T$ and
sums Matsubara modes $\omega_k^t = 2\pi k / \beta$; Wick rotation
analytically continues a Euclidean computation to Lorentzian
signature without changing the periodicity of the underlying time
circle.
Second, target physics:\ observation / temporal-window targets
finite-temperature observables (Casimir at $T \neq 0$, partition
functions, Stefan--Boltzmann); Wick rotation targets
Lorentzian-signature observables read from Euclidean computations
(Hawking radiation, Unruh thermal bath, BW modular flow).
Third, falsifier sharing:\ the operator-level falsifier for
\S~\ref{sec:proj_observation}'s Bisognano--Wichmann footnote
(closing $\sigma_{i\cdot 2\pi}$ on $\mathcal{T}_{n_{\max}}$) is
identical to the falsifier for the four faces of this projection,
because the bridge in every case runs through the same
KMS--Tomita--Takesaki algebra; one operator-system-level proof
lifts all four faces simultaneously.  Same structural correspondence,
same M1 generator, same falsifier, but distinct mechanisms and
distinct target classes of physical observable.  These are not the
same projection.

Cross-references:\ Paper~32~§VIII Remark on the Bisognano--Wichmann
reading documents the spectral-triple-side statement and the
case-exhaustion-theorem context; Paper~35~§VIII subsection on the
Bisognano--Wichmann reading of the temporal-window projection
develops the cigar identification at Hartle--Hawking
$\beta = 8\pi M$; Paper~35~§VIII subsection on the Unruh pendant
develops the Rindler-wedge identification at $\beta = 2\pi / a$;
\S~\ref{sec:proj_observation} carries the original
Bisognano--Wichmann reading footnote on the temporal-window
projection.  The structural ground for the present \S~\ref{sec:proj_wick_rotation}
entry is the Sprint TD Track~4 + Sprint Unruh-pendant + Track~D
documentation chain (May~2026), which promotes the metric-functional-level
content of the Wick-rotation reading from a footnote (\S~\ref{sec:proj_observation})
and an open-question candidate (\S~\ref{sec:open} entry
on the Lorentz boost; see \S~\ref{sec:open}) to a named projection
slot in \S~\ref{sec:layer2}.  The full operator-system-level
Lorentzian extension remains open at the multi-month NCG-mathematics
scale and is documented as a deepening of this entry, not as a
separate candidate.
